\documentclass[10pt,a4paper]{amsart}
\usepackage[margin=19mm]{geometry}
\usepackage{amsmath,amssymb,mathtools}
\usepackage[scr=rsfs,cal=euler]{mathalfa}
\usepackage{xfrac}
\usepackage[unicode,hidelinks,bookmarks=false]{hyperref}
\newcommand{\ie}{i.e.\ }
\newcommand{\cf}{cf.\ }
\newcommand{\resp}{respectively\ }
\newcommand{\Subsection}{\subsection}
\providecommand{\nobreakdash}{\nobreak\hbox{-}}
\setlength{\emergencystretch}{2em}
\multlinegap0pt
\usepackage{amssymb}
\usepackage{xcolor}
\usepackage{tikz}
\newtheorem{prop}{Proposition}
\title{Complete Mappings of Semigroups}
\author{Jo\~ao Ara\'ujo, Wolfram Bentz, Peter J. Cameron, Kevin Hendrey, Michael Kinyon}
\date{Source: arXiv:2608.25092v1 (CC BY 4.0)}
\begin{document}
\maketitle

\noindent\emph{Source and licence.} Complete Mappings of Semigroups. Version arXiv:2608.25092v1 (CC BY 4.0), distributed by arXiv under the Creative Commons Attribution 4.0 International licence, http://creativecommons.org/licenses/by/4.0/, as recorded in the arXiv licence metadata for this exact version and shown on the abstract page https://arxiv.org/abs/2608.25092v1. Full licence notice: \texttt{licenses/arxiv-2608.25092-CC-BY-4.0.txt}.\par\smallskip

\noindent\textit{Editorial scope.} Counted unit: the article's prop p:unique (source0.tex lines 4266--4269). The article's complete original proof of it is delivered with it (source0.tex lines 4271--4451, one \texttt{proof} environment of 181 lines). No mathematics is written, repaired or re-derived: the statement and the proof are the article's own lines, in the article's own notation, and no retained line is substituted or re-flowed.  No further source-local context is needed: every cross-reference of every retained line resolves inside the two delivered spans.  Nothing was omitted from any retained span, so the package carries no omission record.  No retained line contains a citation, so the wrapper carries no bibliography and no bibliography entry.  No retained line contains a figure environment, an image-inclusion command or drawing code, so no image is delivered and the package declares no asset dependency.  Editorial formatting only, in addition: an AMS \texttt{amsart} preamble that restates the article's own declarations and macros used by the retained lines, namely the environments \texttt{prop} on separate counters, together with the standard packages those lines need.  The retained environments print in this document as Proposition 1 in order of appearance, whereas the article prints the same items with the numbering of its own full document.  The complete native source of the article as distributed in the arXiv e-print for this exact version is bundled unchanged as \texttt{sources/208544/source0.tex}.
\begin{prop}\label{p:unique}
    A finite band has exactly one complete mapping if and only if it is a
    semilattice.
\end{prop}

\begin{proof}
    Since $S$ is a band, 
    the identity is a complete mapping.

    We start by proving that if $S$ is not commutative, then it contains a two-element left-zero or right-zero
    subsemigroup.

    Consider the following two assertions:
    \[
        uv=u\ \text{and}\ vu=v
        \ \Rightarrow\ 
        u=v,
        \tag{30}
    \]
    and
    \[
        uv=v\ \text{and}\ vu=u
        \ \Rightarrow\ 
        u=v.
        \tag{31}
    \]
    Suppose, for a contradiction, that both assertions hold for every
    $u,v\in S$. We claim  that $S$ is commutative.

    Let $x,y\in S$. Apply (31) with
    \[
        u=xy
        \ \text{and}\ 
        v=xyx.
    \]
    Since $S$ is a band, both $xy$ and $xyx$ are
    idempotent. Moreover,
    \[
        uv=(xy)(xyx)=xyxyx=(xy)^2x=xyx=v,
    \]
    and
    \[
        vu=(xyx)(xy)=xyxxy=xyxy=(xy)^2=xy=u.
    \]
    Hence (31) implies that
    \[
        xy=xyx.
        \tag{32}
    \]

    Next apply (30) with
    \[
        u=yx
        \ \text{and}\ 
        v=xyx.
    \]
    We have
    \[
        uv=(yx)(xyx)=yxxyx=yxyx=(yx)^2=yx=u.
    \]
    Also,
    \[
        vu=(xyx)(yx)=xyxyx.
    \]
    Since
    \[
        (xyx)^2=xyxxyx=xyxyx
    \]
    and $xyx$ is idempotent, it follows that
    \[
        vu=xyxyx=(xyx)^2=xyx=v.
    \]
    Therefore (30) implies that
    \[
        yx=xyx.
        \tag{33}
    \]
    Combining (32) and (33), we obtain
    \[
        xy=xyx=yx.
    \]
    Since $x$ and $y$ were arbitrary, $S$ is commutative.

    It follows that if $S$ is non-commutative, then at least one of (30)
    and (31) fails. If (30) fails, then there exist distinct elements
    $a,b\in S$ such that
   $ab=a$
        and $ba=b$. 
    Together with $a^2=a$ and $b^2=b$, these equalities show that
   $uv=u$, 
    for all $u,v\in\{a,b\}$. Thus $\{a,b\}$ is a two-element left-zero
    semigroup.

    If (31) fails, it follows by symmetry that  there exist distinct elements
$a,b\in S$ such that $\{a,b\}$ is a two-element right-zero semigroup.

    In either case, the transposition $(a\ b)$ of $S$  is a non-identity complete mapping.
It is proved that if $S$ has exactly one complete mapping, then $S$ is
commutative, and hence it is a semilattice.

    Conversely, suppose that $S$ is a finite semilattice, and let $\phi$
    be a complete mapping of $S$. We prove that $\phi$ is the identity.

    Define the natural partial order on $S$ by
    \[
        x\leq y
        \ \Leftrightarrow\ 
        xy=x.
    \]


    Let
    \[
        \widehat{\phi}\colon S\rightarrow S,
        \ 
        x\widehat{\phi}=x(x\phi).
    \]
    Since $\phi$ is a complete mapping, $\widehat{\phi}$ is a permutation
    of $S$. For every $x\in S$, we have
    \[
        \bigl(x\widehat{\phi}\bigr)x
        =\bigl(x(x\phi)\bigr)x
        =x(x\phi)
        =x\widehat{\phi},
    \]
    using commutativity and idempotence. Hence
    \[x\widehat{\phi}\leq x
        \tag{34}
    \]
    for every $x\in S$.

    We now use the finiteness of $S$. Fix $x\in S$. Since
    $\widehat{\phi}$ is a permutation of the finite set $S$, the element
    $x$ lies in a finite cycle of $\widehat{\phi}$. Thus, for some
    positive integer $m$,
$x{\widehat{\phi}}^{^{^{m}}}=x$.
    Applying (34) successively gives
    \[
        x
        \geq x\widehat{\phi}
        \geq x\widehat{\phi}^{\,2}
        \geq\cdots
        \geq x\widehat{\phi}^{\,m}
        =x.
    \]
    By antisymmetry of the partial order, every inequality in this chain
    is an equality. In particular,
    \[
        x\widehat{\phi}=x.
    \]
    Since $x$ was arbitrary, $\widehat{\phi}$ is the identity
    permutation. Therefore
    \[
        x(x\phi)=x
    \]
    for every $x\in S$, which, by the definition of the natural order,
    means that
    \[
        x\leq x\phi
        \tag{35}
    \]
    for every $x\in S$.

    Finally, $\phi$ is itself a permutation of the finite set $S$. Fix
    $x\in S$, and choose a positive integer $n$ such that
    \[
        x\phi^n=x.
    \]
    Applying (35), we get
    \[
        x
        \leq x\phi
        \leq x\phi^2
        \leq\cdots
        \leq x\phi^n
        =x.
    \]
    Antisymmetry again shows that all these elements are equal. In
    particular,
    \[
        x\phi=x.
    \]
    
    Thus a finite semilattice has exactly one complete mapping, namely
    the identity.
\end{proof}

\end{document}
